\documentclass[11pt]{article}
\usepackage[margin=1.1in]{geometry}
\usepackage{amsmath,amssymb,amsthm}
\usepackage{hyperref}

\newtheorem{theorem}{Theorem}[section]
\newtheorem{lemma}[theorem]{Lemma}
\theoremstyle{definition}
\newtheorem{definition}[theorem]{Definition}
\newtheorem{conjecture}[theorem]{Conjecture}
\theoremstyle{remark}
\newtheorem*{remark}{Remark}

\title{A Machine-Verified Disproof of TLMC Conjecture \#458:\\ The Möbius Number of the Tamari Lattice}
\author{Batch 8 solution submission}
\date{September 12, 2026}

\begin{document}
\maketitle

\begin{abstract}
We disprove a conjecture of the Last Math Competition asserting that the absolute value of the Möbius number of the Tamari lattice $T_n$ is $(n-2)(-1)^n$, ``a closed formula verifiable up to $n = 8$''. In fact the Tamari Möbius numbers are $\mu(T_2) = -1$, $\mu(T_3) = 1$ (the pentagon), $\mu(T_4) = -1$; hence $|\mu(T_4)| = 1 \ne 2 = n - 2$ (and already $|\mu(T_2)| = 1 \ne 0$). The formalization constructs the Tamari lattices $T_2, T_3, T_4$ from scratch (binary trees, right-rotation order), validates the construction against Chapoton's interval-count theorem ($T_3$ has $13$ and $T_4$ has $68$ intervals), and computes the Möbius numbers by dynamic programming of the textbook recursion $\mu(x,y) = -\sum_{x \le z < y} \mu(x,z)$ along a linear extension. All computations are kernel-verified; the main results depend on at most \texttt{[propext]}.
\end{abstract}

\section{The conjecture}

\begin{conjecture}[TLMC-00000000458]\label{conj:458}
The absolute value of the Möbius number of the Tamari lattice $T_n$ is $(n-2)(-1)^n$ (a closed formula verifiable up to $n = 8$).
\end{conjecture}

\section{The disproof}

\begin{definition}[Tamari lattice]
$T_n$ is the lattice of binary trees with $n$ internal nodes, ordered by the reachability of right rotations $((a,b),c) \mapsto (a,(b,c))$. The minimum is the left comb and the maximum the right comb.
\end{definition}

\begin{theorem}\label{thm:458}
Conjecture~\ref{conj:458} is \textbf{false}:
\[
\mu(T_2) = -1, \qquad \mu(T_3) = 1, \qquad \mu(T_4) = -1,
\]
so $|\mu(T_4)| = 1 \ne 2 = n-2$, and also $|\mu(T_2)| = 1 \ne 0$.
\end{theorem}

\begin{proof}
The Möbius function is computed by the textbook recursion $\mu(x,x) = 1$, $\mu(x,y) = -\sum_{x \le z < y} \mu(x,z)$, evaluated by dynamic programming along a linear extension of $T_n$ (the trees sorted by decreasing closure size). For $T_3$ (the pentagon) this gives the classical value $\mu = 1$: the two atoms contribute $-1$ each and the coatom between them $0$. For $T_4$ (14 elements) the recursion yields $-1$, not $\pm 2$.
\end{proof}

\begin{remark}[Validation of the construction]
The poset built here is the genuine Tamari lattice: the numbers of intervals of $T_3$ and $T_4$ come out as $13$ and $68$, matching Chapoton's theorem $\#\text{intervals} = \frac{2}{n(n+1)}\binom{4n+1}{n-1}$ ($T_5$ gives $399$ in the accompanying computation). The min/max comb bounds and the linear-extension property are also verified.
\end{remark}

\begin{remark}
The ``verified up to $n=8$'' claim of the conjecture cannot refer to the standard Tamari lattice: its Möbius numbers do not grow linearly. The formula $(n-2)$ agrees with the truth only at $n = 3$ (the pentagon, $|\mu| = 1$).
\end{remark}

\section{Machine verification}

\texttt{Batch8.lean} (namespace \texttt{TLMCBatch8}) contains: \texttt{trees}/\texttt{treesA} (fuel-structural enumeration of binary trees, with \texttt{trees\_stable}), \texttt{coversOf} (right rotations via a rewrite system), \texttt{closureN}/\texttt{leT} (reachability, with \texttt{closure\_stable}), \texttt{interval\_counts} (Chapoton validation), \texttt{comb\_bounds}, \texttt{extOf}/\texttt{ext\_linear}, the Möbius DP \texttt{muRow}/\texttt{muTop}, and the refutation \texttt{tlm458\_false} with the packaged data \texttt{tlm458}. Compilation via \texttt{lake env lean Batch8.lean} exits $0$ with zero errors and warnings; \texttt{\#print axioms} shows \texttt{trees\_cards}, \texttt{trees\_stable}, \texttt{tlm458\_false} depend on \emph{no} axioms, and the remaining results on at most \texttt{[propext]}. No \texttt{sorry} appears.

\end{document}
